% =============================================================================
% Shared preamble of the real-exercise list for Lectures 1--5 (real schedule).
% Input by Listas/lista-reais-lectures-01-05.tex (all blocks) and by
% Listas/lista-reais-lecture-0N.tex (one block). The driver defines \listatema before
% \input-ing this file; a per-lecture driver then sets \perlecturetrue and \blockoffset
% (number of items in the earlier blocks) so items are numbered 1--10 and also cite their
% number in the combined list. Rules: rules/00_estilo_avaliacao.md §5.
% =============================================================================
\usepackage[utf8]{inputenc}
\usepackage[T1]{fontenc}
\usepackage{lmodern}   % Type 1 vector fonts (replaces PK bitmaps)
\usepackage{cmap}      % ToUnicode CMap -> copyable text and accents
\usepackage[english]{babel}
\usepackage{amsmath,amssymb,amsfonts,mathtools}
\usepackage[a4paper,margin=2.2cm]{geometry}
\usepackage{enumitem}
\usepackage{fancyhdr}
\usepackage{xcolor}
\usepackage{microtype}
\usepackage{booktabs}
\usepackage{graphicx}
\usepackage{array}
\usepackage{needspace}

\newcommand{\disciplina}{Economia das Organizações e dos Incentivos --- MPE Insper}
\newcommand{\espacoresol}{3.5cm}   % height of the resolution space after each exercise

\pagestyle{fancy}\fancyhf{}
\lhead{\footnotesize\disciplina}
\rhead{\footnotesize \listatema}
\cfoot{\footnotesize\thepage}
\renewcommand{\headrulewidth}{0.3pt}
\setlength{\headheight}{14pt}

\setlist{topsep=2pt,itemsep=2pt}
\setlength{\parindent}{0pt}
\setlength{\parskip}{4pt}
\newcounter{questao}

% Per-lecture mode: items numbered 1--10, with their number in the combined list.
\newif\ifperlecture
\newcommand{\blockoffset}{0}
\newcommand{\rownum}[1]{\ifperlecture\the\numexpr#1-\blockoffset\relax\ (#1)\else#1\fi}

% Block header: number, lecture title, coverage.
\newcommand{\bloco}[3]{\bigskip\noindent
  {\large\bfseries\color{blue!55!black}Block #1 --- #2}\hfill\textit{(#3)}\par
  \noindent\rule{\linewidth}{0.3pt}\smallskip}
% Exercise header: [#1] difficulty 1--5, #2 short name, #3 level, #4 source line.
\newcommand{\exercicio}[4][]{\needspace{12\baselineskip}\refstepcounter{questao}\medskip\noindent
  \textbf{Exercise \thequestao.}\ifperlecture~{\small(item
  \the\numexpr\value{questao}+\blockoffset\relax\ of the combined list)}\fi~\textit{#2}%
  % keeps the [level; difficulty] tag whole, right-aligned on the last line (TeXbook \signed)
  {\unskip\nobreak\hfil\penalty50\hskip1em\hbox{}\nobreak\hfill\mbox{\small[#3; difficulty #1/5]}}\par
  \noindent{\small\color{black!70}\textbf{Source.} #4}\par\smallskip}
% Material an exercise refers to, copied in under its own source line.
\newenvironment{referencia}[1]{\par\smallskip\noindent
  \begin{minipage}{\linewidth}\small\leftskip=1.2em
  \noindent\textbf{Referenced material.} {\color{black!70}\textbf{Source.} #1}\par\smallskip}%
  {\end{minipage}\par}
% Book's exercise label inside the statement, e.g. "4.5".
\newcommand{\rot}[1]{\textbf{#1}\quad}
\newcommand{\resolespaco}{\par\textit{Resolution:}\vspace{\espacoresol}\par}
% Payoff matrix: #1 row-player label, #2 column-player label, #3 inner tabular.
\newcommand{\jogo}[3]{\begin{center}\begin{tabular}{@{}r@{\hspace{6pt}}c@{}}
  & #2\\[2pt] #1 & #3\end{tabular}\end{center}}
\newcommand{\hd}[1]{\multicolumn{1}{c}{#1}}
\newenvironment{partes}{\begin{enumerate}[label=\alph*.,leftmargin=2em]}{\end{enumerate}}

% ----------------------------------------------------------------------------- front page
\newcommand{\titulo}{\begin{center}
  {\Large\bfseries \disciplina}\\[2pt]
  {\large Exercise List: \listatema}
\end{center}}

\newcommand{\howtouse}{{\small\noindent\textbf{How to use this list.} Every exercise below is a
real exercise from one of the course books, reproduced verbatim: the source line gives author,
book, chapter, exercise, the printed page and, in brackets, the page your PDF viewer shows for the
file in \texttt{Livros/}, and whether the source gives the answer. Muñoz-Garcia \& Toro-Gonzalez
print a worked answer right after each exercise; MWG answers are in the MWG solutions manual in
\texttt{Livros/}; Tadelis, Gibbons and Cabral give no answer in the book. When an exercise refers
to a model, game or figure stated elsewhere in its book, that material is copied into the same
block under ``Referenced material'', so every exercise can be solved where it is printed. There is
one block per lecture of the real schedule. Each block has ten exercises in strictly rising
difficulty, rated from 1 (apply a definition) to 5 (graduate level) next to each heading: three
foundations, five exam level, two advanced. Try each exercise in the space provided before opening
the book's answer.\par}}

% Selection table; #1 = the rows of one block (\rowsi ... \rowsv).
\newcommand{\selectiontable}[1]{\begin{center}
\footnotesize
\setlength{\tabcolsep}{2.8pt}
\begin{tabular}{@{}c r l r l >{\raggedright\arraybackslash}p{4.0cm} c c@{}}
\toprule
Block & \# & Source & p. [PDF] & Level & What it trains & Diff. & Answer \\
\midrule
#1
\bottomrule
\end{tabular}
\end{center}}

\newcommand{\tablenote}{{\footnotesize \ifperlecture \#: item number in this lecture's list,
then its number in the combined list (\texttt{lista-reais-lectures-01-05.pdf}) in brackets.
\fi Items not in the list for Lectures 1--4 (\texttt{lista-reais-lectures-01-04.tex}) are marked:
$^\dagger$\,added in the 30-item version, $^\ddagger$\,added in the 50-item version. Every other
item is reused from that list with its statement unchanged, some in a different block or position.
``Diff.'': difficulty from 1 (apply a definition) to 5 (graduate level). Pages: printed page, then
the PDF viewer page in brackets (offsets in \texttt{Livros/book-registry.json}; MWG is a
two-pages-per-sheet scan). ``Answer'': \emph{book} = worked answer printed after the exercise;
\emph{manual} = MWG solutions manual in \texttt{Livros/}; \emph{none} = no answer in the book.
Books: Tadelis (2013), \emph{Game Theory: An Introduction}; Gibbons (1992), \emph{Game Theory for
Applied Economists}; Muñoz-Garcia \& Toro-Gonzalez (2019), \emph{Strategy and Game Theory: Practice
Exercises with Answers}, 2nd ed.; Mas-Colell, Whinston \& Green (1995), \emph{Microeconomic Theory}
(MWG); Cabral (2017), \emph{Introduction to Industrial Organization}, 2nd ed.\par}}

\newcommand{\rowsi}{%
L1 & \rownum{1} & Tadelis 4.5 & 77 [96] & Found. & IESDS round by round, $3\times3$ matrix & 1 & none \\
   & \rownum{2} & Muñoz-Garcia Ch.~1 Ex.~13$^\ddagger$ & 26 [39] & Found. & Dominant and dominated strategies, then IDSDS & 1 & book \\
   & \rownum{3} & Gibbons 1.2 & 48 [56] & Found. & IESDS, then pure Nash equilibria & 2 & none \\
   & \rownum{4} & Tadelis 4.1$^\ddagger$ & 76 [95] & Exam & Proof: a strictly dominant-strategy equilibrium is unique & 2 & none \\
   & \rownum{5} & Muñoz-Garcia Ch.~1 Ex.~12 & 23 [36] & Exam & Rationality vs.\ common knowledge of it & 2 & book \\
   & \rownum{6} & Tadelis 4.7 & 77--78 [96--97] & Exam & Story to matrix, then elimination & 3 & none \\
   & \rownum{7} & Tadelis 4.3$^\ddagger$ & 76 [95] & Exam & Discrete first-price auction: matrix, IESDS & 3 & none \\
   & \rownum{8} & Tadelis 4.6 & 77 [96] & Exam & IESDS with continuous strategies & 4 & none \\
   & \rownum{9} & MWG 8.B.4$^\ddagger$ & 262 [142] & Adv. & Proof: the order of strict deletion does not matter & 4 & manual \\
   & \rownum{10} & MWG 8.B.5 & 262 [142] & Adv. & Cournot by IESDS: two firms vs.\ three & 5 & manual \\}
\newcommand{\rowsii}{%
L2 & \rownum{11} & Tadelis 5.6 & 96 [115] & Found. & Hawk-dove: matrix, pure equilibria & 1 & none \\
   & \rownum{12} & Muñoz-Garcia Ch.~2 Ex.~2$^\ddagger$ & 41 [54] & Found. & Battle of the sexes: pure equilibria & 1 & book \\
   & \rownum{13} & Muñoz-Garcia Ch.~3 Ex.~1 & 86--87 [99--100] & Found. & Chicken: pure and mixed equilibria & 2 & book \\
   & \rownum{14} & Gibbons 1.12$^\ddagger$ & 51 [59] & Exam & A mixed equilibrium by indifference & 2 & none \\
   & \rownum{15} & Tadelis 5.8 & 96--97 [115--116] & Exam & Cournot with $n$ firms; competitive limit & 3 & none \\
   & \rownum{16} & Gibbons 1.6 & 49 [57] & Exam & Cournot, asymmetric costs; corner case & 3 & none \\
   & \rownum{17} & Gibbons 1.7$^\ddagger$ & 49--50 [57--58] & Exam & Homogeneous Bertrand: $p_1 = p_2 = c$ is the unique equilibrium & 3 & none \\
   & \rownum{18} & Tadelis 6.8 & 125 [144] & Exam & Three-player entry; symmetric mixing & 4 & none \\
   & \rownum{19} & MWG 8.D.2$^\ddagger$ & 263 [142] & Adv. & Proof: a unique IESDS survivor is a Nash equilibrium & 4 & manual \\
   & \rownum{20} & MWG 8.D.6 & 264 [143] & Adv. & A mixed equilibrium exists in every $2\times2$ game & 5 & manual \\}
\newcommand{\rowsiii}{%
L3 & \rownum{21} & Tadelis 7.1$^\ddagger$ & 147 [166] & Found. & Counting pure strategies over information sets & 1 & none \\
   & \rownum{22} & Muñoz-Garcia Ch.~4 Ex.~3 & 160 [172] & Found. & Sequential prisoner's dilemma & 1 & book \\
   & \rownum{23} & Tadelis 7.2$^\dagger$ & 147 [166] & Found. & Story to tree; strategies; an ``unappealing'' Nash equilibrium & 2 & none \\
   & \rownum{24} & Tadelis 8.5$^\dagger$ & 170 [189] & Exam & A simultaneous game behind an adoption choice: SPE & 3 & none \\
   & \rownum{25} & Tadelis 8.10 & 172 [191] & Exam & Subgames with a simultaneous stage & 3 & none \\
   & \rownum{26} & Tadelis 7.4$^\ddagger$ & 148 [167] & Exam & Four-move centipede: strategies, Nash equilibria, outcomes & 3 & none \\
   & \rownum{27} & Tadelis 7.9$^\ddagger$ & 149--150 [168--169] & Exam & A move of Nature and an information set & 4 & none \\
   & \rownum{28} & Tadelis 7.10$^\dagger$ & 150 [169] & Exam & Mixed vs.\ behavioral strategies; perfect recall & 4 & none \\
   & \rownum{29} & MWG 9.B.4$^\ddagger$ & 302 [162] & Adv. & Proof: unique SPNE payoffs in zero-sum games of perfect information & 4 & manual \\
   & \rownum{30} & Gibbons 2.1 & 130 [137] & Adv. & Rotten Kid theorem (continuous actions) & 5 & none \\}
\newcommand{\rowsiv}{%
L4 & \rownum{31} & Gibbons 2.12$^\dagger$ & 135 [142] & Found. & Strategy, subgame and SPE in a repeated game & 1 & none \\
   & \rownum{32} & Tadelis 7.6$^\ddagger$ & 148--149 [167--168] & Found. & Entry game: payoffs from preferences, all Nash equilibria & 1 & none \\
   & \rownum{33} & Tadelis 8.8 & 171--172 [190--191] & Found. & Entry deterrence: a non-credible threat & 2 & none \\
   & \rownum{34} & Tadelis 8.6$^\dagger$ & 170--171 [189--190] & Exam & Sunk investment before Cournot; threshold $F^*$ & 3 & none \\
   & \rownum{35} & Muñoz-Garcia Ch.~4 Ex.~11$^\ddagger$ & 180 [192] & Exam & Stackelberg, asymmetric costs; vs.\ monopoly, Cournot, Bertrand & 3 & book \\
   & \rownum{36} & Tadelis 8.11$^\ddagger$ & 172--173 [191--192] & Exam & Sequential vetoes; Nash equilibria that are not SPE; removing an option & 3 & none \\
   & \rownum{37} & Tadelis 8.4$^\dagger$ & 170 [189] & Exam & Leader, two followers: subgames; a Nash equilibrium that is not SPE & 4 & none \\
   & \rownum{38} & Tadelis 8.13 & 173--174 [192--193] & Exam & Hold-up; a posted price as commitment & 4 & none \\
   & \rownum{39} & Tadelis 8.7$^\ddagger$ & 171 [190] & Adv. & Credibility of suing; design of the legal system & 4 & none \\
   & \rownum{40} & MWG 9.B.5$^\dagger$ & 302 [162] & Adv. & Observation, multiple SPNE, when moving first pays & 5 & manual \\}
\newcommand{\rowsv}{%
L5 & \rownum{41} & Cabral 7.13$^\dagger$ & 182 [194] & Found. & Finitely repeated PD unravels & 1 & none \\
   & \rownum{42} & Tadelis 10.2$^\dagger$ & 215 [234] & Found. & Grim trigger; threshold $\delta$ for two paths & 2 & none \\
   & \rownum{43} & Muñoz-Garcia Ch.~6 Ex.~1$^\ddagger$ & 274 [285] & Found. & One-, two- and ten-period punishments & 2 & book \\
   & \rownum{44} & Muñoz-Garcia Ch.~6 Ex.~9$^\dagger$ & 305 [316] & Exam & Feasible, individually rational payoffs; folk theorem & 3 & book \\
   & \rownum{45} & Tadelis 10.3$^\ddagger$ & 215--216 [234--235] & Exam & Length-$T$ punishments: $\delta_1$ vs.\ $\delta_2$ & 3 & none \\
   & \rownum{46} & Tadelis 9.1 & 186--187 [205--206] & Exam & Stage game played twice: all SPE, carrot and stick & 4 & none \\
   & \rownum{47} & Gibbons 2.11$^\ddagger$ & 135 [142] & Exam & A non-Nash first-stage payoff in a twice-played game & 4 & none \\
   & \rownum{48} & Tadelis 11.5 & 236--237 [255--256] & Exam & Three periods with $\delta_1\neq\delta_2$ & 4 & none \\
   & \rownum{49} & Tadelis 10.7$^\ddagger$ & 217--218 [236--237] & Adv. & Quality: ten days vs.\ forever, in parameters & 4 & none \\
   & \rownum{50} & Gibbons 2.20$^\dagger$ & 137--138 [144--145] & Adv. & Rubinstein strategies are an SPE (one-stage deviations) & 5 & none \\}

% Difficulty path: what makes each block's last item hard.
\newcommand{\pathi}{\textbf{Lecture 1}: from one $3\times3$ elimination (1) through two short proofs
about dominance (Tadelis 4.1, MWG 8.B.4) to MWG 8.B.5 (5), where elimination over a continuum of
quantities must be shown to converge with two firms and to stall with three.}
\newcommand{\pathii}{\textbf{Lecture 2}: from writing a matrix (1) through the Bertrand uniqueness
proof and MWG 8.D.2 to MWG 8.D.6 (5), a proof that every $2\times2$ game, with arbitrary payoffs,
has a mixed-strategy equilibrium.}
\newcommand{\pathiii}{\textbf{Lecture 3}: from counting strategies (1) through Nature, imperfect
recall and the zero-sum proof (MWG 9.B.4) to Gibbons 2.1 (5), a proof with general concave income
and utility functions and continuous actions.}
\newcommand{\pathiv}{\textbf{Lecture 4}: from definitions (1) through entry, Stackelberg, hold-up and
the legal-system game (Tadelis 8.7) to MWG 9.B.5 (5), a five-part general-payoff argument on
strategy counts, multiple SPNE and when moving first guarantees the simultaneous-move payoff, with
counterexamples.}
\newcommand{\pathv}{\textbf{Lecture 5}: from unraveling (1) through punishments of finite length,
carrot and stick and the both-horizons bar game (Tadelis 10.7) to Gibbons 2.20 (5), a proof that
Rubinstein's stationary strategies are a Nash equilibrium and are subgame-perfect, checking every
history of an infinite game.}
